\documentclass{article} 
\input{../../../lib.sty} 
 
\title{\texttt{fn sample\_uniform\_uint\_below}} 
\author{\MichaelShoemateAuthor} 
 
\begin{document} 
\maketitle 
 
This document proves that the implementation of \rustdoc{traits/samplers/uniform/fn}{sample\_uniform\_uint\_below} in \asOfCommit{mod.rs}{f5bb719} 
satisfies its definition. 
 
\texttt{sample\_uniform\_uint\_below} uses rejection sampling. 
In each round all bits of the integer are filled randomly, drawing an unsigned integer uniformly at random from a space of size $2^N$. 
The algorithm returns the sample, modulo the upper bound, so long as the sample is not one of the first $2^N \mod \texttt{upper}$ integers. 
 
\section{Hoare Triple} 
\subsection*{Precondition} 
 
\subsubsection*{Compiler-verified} 
\begin{itemize} 
    \item Generic \texttt{T} has traits \texttt{Integer + Unsigned + FromBytes<N>}. 
    \item Const-generic \texttt{N} of type usize. By the definition of \texttt{FromBytes}, the compiler ensures this is the number of bytes in T. 
    \item Argument \texttt{upper} must be of type \texttt{T} 
\end{itemize} 
 
\subsubsection*{Caller-verified} 
$\texttt{upper} > 0$ 
 
\subsection*{Pseudocode} 
 
\lstinputlisting[language=Python,firstline=2]{./pseudocode/sample_uniform_uint_below.py} 
 
\subsection*{Postcondition} 
For any setting of the input parameter \texttt{upper} that satisfies the precondition, 
\texttt{sample\_uniform\_uint\_below} either 
\begin{itemize} 
    \item raises an exception if there is a lack of system entropy, 
    \item returns \texttt{out} where \texttt{out} is uniformly distributed between $[0, upper)$. 
\end{itemize} 
 
\begin{proof} 
\label{unsigned-integer-proof} 
By the postcondition of \rustdoc{traits/samplers/uniform/fn}{sample\_from\_uniform\_bytes}, 
\texttt{sample} is drawn uniformly from the $2^N$ representable values of type \texttt{T}, namely $\{0, \ldots, \texttt{T.MAX}\}$. 
 
Let $r = 2^N \mod \texttt{upper}$. Since \texttt{upper} is positive by the preconditions, $r$ is well-defined and $\texttt{reject\_below} = r$. 
The algorithm rejects exactly the first $r$ values and accepts the remaining $2^N - r$ values. 
This is equivalent to the more familiar presentation that accepts the first $2^N - r$ values and rejects the final $r$ values, because both acceptance regions have the same size and therefore contain the same number of representatives from each congruence class modulo \texttt{upper}. 
The implementation uses front rejection because $2^N = \texttt{T.MAX} + 1$ may not itself be representable in type \texttt{T}$;$ however, $r$ is representable, so testing whether \texttt{sample >= reject\_below} avoids overflow. 
Because $2^N - r$ is a multiple of \texttt{upper}, the accepted values partition into congruence classes modulo \texttt{upper} of equal size. 
Therefore \texttt{sample \% upper} is uniformly distributed on $[0, \texttt{upper})$ after conditioning on acceptance. 
 
\noindent Therefore, for any positive value of \texttt{upper}, the function satisfies the postcondition. 
\end{proof} 
 
\end{document} 
